% GENERATED FILE. Edit canonical lessons, then run npm run generate:books.
% canonical-source-hash: fnv1a64:7c1e63c82c677905
% canonical-lessons: HI-C313-patrakar, HI-C313-ghatnasthal, HI-C313-chaapna, HI-C313-sampadak, HI-C313-vivaran

\chapter{Reporters}
\label{ch:hi-reporters}
\hlchaptermodality{313}

\begin{chapteropening}
\textbf{By the end of this chapter:} \emph{I can say who reported a story and from where.}

Complete the last lesson of chapter 313: I can say who reported a story and from where.
\end{chapteropening}

\section[\texorpdfstring{\hi{पत्रकार} (patrakār)}{पत्रकार (patrakār)}]{\hi{पत्रकार} (patrakār) — a journalist}

\label{lesson:HI-C313-patrakar}



\begin{quote}
Before the new one: say the Hindi for a headline, then the Hindi for to be printed.
\end{quote}

\subsection*{You'll want to know: \hi{पत्रकार}}
\textbf{\hi{पत्रकार}} — \emph{patrakār} — \textquotedblleft{}a journalist\textquotedblright{}.

\begin{quote}
\hi{मीरा} \hi{पत्रकार} \hi{बनना} \hi{चाहती} \hi{है।}
\end{quote}

\begin{quote}
\emph{Mīrā patrakār bannā cāhtī hai.}
\end{quote}

\begin{quote}
Mira wants to become a journalist.
\end{quote}

Masculine or feminine, by the person meant.

\begin{cousinweb}
\textbf{\hi{पत्रकार}} (\emph{patrakār}) joins \textbf{\hi{पत्र}}, a letter or a paper, from Sanskrit \emph{pattra}, also spelt \emph{patra}, with \emph{-kāra}, \textquotedblleft{}maker, doer\textquotedblright{} (Monier-Williams): a modern word for one who makes the papers, a journalist.
\end{cousinweb}

\begin{grammarlens}[title={What you've built}]
One more word for telling what happened.
\end{grammarlens}

\subsection*{Your turn}
\begin{itemize}
\raggedright
  \item \emph{Say it:} \emph{patrakār}
  \item \emph{Say it:} \emph{patrakār}, once more
  \item \emph{Say it:} \emph{Mīrā patrakār bannā cāhtī hai.} — Mira wants to become a journalist.
  \item \emph{Recall:} say the Hindi for a headline, then the Hindi for to be printed, then say \emph{patrakār} again
\end{itemize}

\begin{tcolorbox}[breakable,colback=teal!4,colframe=teal!35!black,title={Before you move on}]
What does \hi{पत्रकार} mean? (\textquotedblleft{}A journalist\textquotedblright{}.) Say it once more.
\end{tcolorbox}
\section[\texorpdfstring{\hi{घटनास्थल} (ghaṭnāsthal)}{घटनास्थल (ghaṭnāsthal)}]{\hi{घटनास्थल} (ghaṭnāsthal) — the scene (of an incident)}

\label{lesson:HI-C313-ghatnasthal}



\begin{quote}
Before the new one: say the Hindi for to be printed, then the Hindi for a journalist.
\end{quote}

\subsection*{You'll want to know: \hi{घटनास्थल}}
\textbf{\hi{घटनास्थल}} — \emph{ghaṭnāsthal} — \textquotedblleft{}the scene (of an incident)\textquotedblright{}.

\begin{quote}
\hi{घटनास्थल} \hi{के} \hi{लिए} \hi{पुल} \hi{के} \hi{बाद} \hi{बाएँ} \hi{मुड़िए।}
\end{quote}

\begin{quote}
\emph{ghaṭnāsthal ke liye pul ke bād bāẽ muṛiye.}
\end{quote}

\begin{quote}
For the scene of the accident, turn left after the bridge.
\end{quote}

Masculine.

\begin{cousinweb}
\textbf{\hi{घटनास्थल}} (\emph{ghaṭnāsthal}) joins \textbf{\hi{घटना}}, an event, with Sanskrit \emph{sthala}, \textquotedblleft{}a place\textquotedblright{} (Monier-Williams): the place where it happened. It is the word of news reports.
\end{cousinweb}

\begin{grammarlens}[title={What you've built}]
One more word for telling what happened.
\end{grammarlens}

\subsection*{Your turn}
\begin{itemize}
\raggedright
  \item \emph{Say it:} \emph{ghaṭnāsthal}
  \item \emph{Say it:} \emph{ghaṭnāsthal}, once more
  \item \emph{Say it:} \emph{ghaṭnāsthal ke liye pul ke bād bāẽ muṛiye.} — For the scene of the accident, turn left after the bridge.
  \item \emph{Recall:} say the Hindi for to be printed, then the Hindi for a journalist, then say \emph{ghaṭnāsthal} again
\end{itemize}

\begin{tcolorbox}[breakable,colback=teal!4,colframe=teal!35!black,title={Before you move on}]
What does \hi{घटनास्थल} mean? (\textquotedblleft{}The scene (of an incident)\textquotedblright{}.) Say it once more.
\end{tcolorbox}
\section[\texorpdfstring{\hi{छापना} (chāpnā)}{छापना (chāpnā)}]{\hi{छापना} (chāpnā) — to print}

\label{lesson:HI-C313-chaapna}



\begin{quote}
Before the new one: say the Hindi for a journalist, then the Hindi for the scene (of an incident).
\end{quote}

\subsection*{You'll want to know: \hi{छापना}}
\textbf{\hi{छापना}} — \emph{chāpnā} — \textquotedblleft{}to print\textquotedblright{}.

\begin{quote}
\hi{अख़बार} \hi{ने} \hi{हादसे} \hi{की} \hi{तस्वीर} \hi{छापी।}
\end{quote}

\begin{quote}
\emph{akhbār ne hādse kī tasvīr chāpī.}
\end{quote}

\begin{quote}
The newspaper printed a picture of the accident.
\end{quote}

Past, with \hi{ने} (he or she …): \hi{उसने} \hi{ख़बर} \hi{छापी}, \emph{usne khabar chāpī}.

\begin{cousinweb}
\textbf{\hi{छापना}} (\emph{chāpnā}) is inherited from \emph{*chapp}, a pressing, a form CDIAL reconstructs (CDIAL 4994): printing first meant pressing a mark on to cloth or paper.
\end{cousinweb}

\begin{grammarlens}[title={What you've built}]
One more word for telling what happened.
\end{grammarlens}

\subsection*{Your turn}
\begin{itemize}
\raggedright
  \item \emph{Say it:} \emph{chāpnā}
  \item \emph{Say it:} \emph{chāpnā}, once more
  \item \emph{Say it:} \emph{akhbār ne hādse kī tasvīr chāpī.} — The newspaper printed a picture of the accident.
  \item \emph{Recall:} say the Hindi for a journalist, then the Hindi for the scene (of an incident), then say \emph{chāpnā} again
\end{itemize}

\begin{tcolorbox}[breakable,colback=teal!4,colframe=teal!35!black,title={Before you move on}]
What does \hi{छापना} mean? (\textquotedblleft{}To print\textquotedblright{}.) Say it once more.
\end{tcolorbox}
\section[\texorpdfstring{\hi{संपादक} (sampādak)}{संपादक (sampādak)}]{\hi{संपादक} (sampādak) — an editor}

\label{lesson:HI-C313-sampadak}



\begin{quote}
Before the new one: say the Hindi for the scene (of an incident), then the Hindi for to print.
\end{quote}

\subsection*{You'll want to know: \hi{संपादक}}
\textbf{\hi{संपादक}} — \emph{sampādak} — \textquotedblleft{}an editor\textquotedblright{}.

\begin{quote}
\hi{संपादक} \hi{ने} \hi{ख़बर} \hi{को} \hi{दो} \hi{बार} \hi{पढ़ा।}
\end{quote}

\begin{quote}
\emph{sampādak ne khabar ko do bār paṛhā.}
\end{quote}

\begin{quote}
The editor read the story twice.
\end{quote}

Masculine.

\begin{cousinweb}
\textbf{\hi{संपादक}} (\emph{sampādak}) is Sanskrit \emph{sampādaka}, \textquotedblleft{}one who brings about, who accomplishes\textquotedblright{} (Monier-Williams). The sense of editor is modern: the one who gets the paper ready.
\end{cousinweb}

\begin{grammarlens}[title={What you've built}]
One more word for telling what happened.
\end{grammarlens}

\subsection*{Your turn}
\begin{itemize}
\raggedright
  \item \emph{Say it:} \emph{sampādak}
  \item \emph{Say it:} \emph{sampādak}, once more
  \item \emph{Say it:} \emph{sampādak ne khabar ko do bār paṛhā.} — The editor read the story twice.
  \item \emph{Recall:} say the Hindi for the scene (of an incident), then the Hindi for to print, then say \emph{sampādak} again
\end{itemize}

\begin{tcolorbox}[breakable,colback=teal!4,colframe=teal!35!black,title={Before you move on}]
What does \hi{संपादक} mean? (\textquotedblleft{}An editor\textquotedblright{}.) Say it once more.
\end{tcolorbox}
\section[\texorpdfstring{\hi{विवरण} (vivaraṇ)}{विवरण (vivaraṇ)}]{\hi{विवरण} (vivaraṇ) — a detailed account}

\label{lesson:HI-C313-vivaran}



\begin{quote}
Before the new one: say the Hindi for to print, then the Hindi for an editor.
\end{quote}

\subsection*{You'll want to know: \hi{विवरण}}
\textbf{\hi{विवरण}} — \emph{vivaraṇ} — \textquotedblleft{}a detailed account\textquotedblright{}.

\begin{quote}
\hi{पत्रकार} \hi{ने} \hi{हादसे} \hi{का} \hi{पूरा} \hi{विवरण} \hi{दिया।}
\end{quote}

\begin{quote}
\emph{patrakār ne hādse kā pūrā vivaraṇ diyā.}
\end{quote}

\begin{quote}
The reporter gave a full account of the accident.
\end{quote}

Masculine.

\begin{cousinweb}
\textbf{\hi{विवरण}} (\emph{vivaraṇ}) is Sanskrit \emph{vivaraṇa}, \textquotedblleft{}an opening up, an explaining\textquotedblright{} (Monier-Williams): the facts laid open one by one.
\end{cousinweb}

\begin{grammarlens}[title={What you've built}]
One more word for telling what happened.
\end{grammarlens}

\subsection*{Your turn}
\begin{itemize}
\raggedright
  \item \emph{Say it:} \emph{vivaraṇ}
  \item \emph{Say it:} \emph{vivaraṇ}, once more
  \item \emph{Say it:} \emph{patrakār ne hādse kā pūrā vivaraṇ diyā.} — The reporter gave a full account of the accident.
  \item \emph{Recall:} say the Hindi for to print, then the Hindi for an editor, then say \emph{vivaraṇ} again
\end{itemize}

\begin{tcolorbox}[breakable,colback=teal!4,colframe=teal!35!black,title={Before you move on}]
What does \hi{विवरण} mean? (\textquotedblleft{}A detailed account\textquotedblright{}.) Say it once more.
\end{tcolorbox}
